\section{Методы предобработки и извлечения признаков электроэнцефалографических сигналов}

\subsection{Предобработка электроэнцефалографических сигналов}

Качество последующего декодирования EEG-сигналов в значительной степени определяется корректностью этапов предварительной обработки (рис. \ref{fig:eeg_preprocessing_pipeline}). Низкое отношение сигнал/шум (signal-to-noise ratio, SNR), влияние миографических и окулярных артефактов, а также вариативность условий регистрации делают этап предобработки критически важным для задач машинного обучения. Современные обзоры подчёркивают, что различия в пайплайнах предобработки могут существенно влиять на итоговую точность классификации и воспроизводимость результатов~\cite{10.1109/access.2024.3360328}.

\subsubsection*{Реферирование и фильтрация}

Первым этапом обработки является выбор схемы реферирования (referencing). Поскольку измеряемый потенциал всегда определяется относительно опорного электрода, некорректный выбор референса может привести к искажению пространственной структуры сигнала. Распространённой практикой является использование среднего референса (average reference), однако наличие шумных или «плохих» каналов может существенно исказить усреднение. В этой связи предложены робастные процедуры реферирования, учитывающие автоматическое выявление и исключение некорректных каналов~\cite{10.3389/fninf.2015.00016}.

Полосовая фильтрация (band-pass filtering) применяется для удаления медленных дрейфов (low-frequency drift) и высокочастотного шума, включая сетевые помехи. Типичный диапазон для когнитивных задач составляет 0.1-40 Гц или 1-40 Гц, однако выбор параметров фильтра зависит от исследуемых частотных диапазонов. Важно учитывать, что чрезмерно агрессивная фильтрация может искажать временную структуру сигнала и латентности компонентов.

% КАРТИНКА 1
% prompt: Basic preprocessing pipeline diagram: Raw EEG → Bad channel detection → Re-referencing → Filtering → Artifact correction → Segmentation(**language: Russian**; show signal view on each step; without duplicating the text)
\begin{figure}[!t]
  \centering
  \includegraphics[width=150mm]{eeg_preprocessing_pipeline.png}
  \caption{Схема конвейера предобработки ЭЭГ-данных, включающего последовательные этапы обнаружения и исключения плохих каналов, перереференсирования, частотной фильтрации, коррекции артефактов и сегментации непрерывной записи на отдельные эпохи.}
    \label{fig:eeg_preprocessing_pipeline}
\end{figure}

\subsubsection*{Удаление артефактов}

Одной из ключевых задач предобработки является удаление артефактов, обусловленных миганием, движениями глаз, мышечной активностью и электродными смещениями. Широко применяемым методом является независимый компонентный анализ (Independent Component Analysis, ICA), позволяющий разложить многоканальный сигнал на статистически независимые компоненты и выделить артефактные источники~\cite{10.1016/j.jneumeth.2003.10.009}. 

Использование инструментов, таких как программный пакет EEGLAB, значительно упростило применение ICA и сделало возможным анализ сигналов на уровне отдельных проб (single-trial analysis). Однако эффективность ICA зависит от качества исходных данных и предварительной фильтрации, а также от корректной идентификации компонент.

\subsubsection*{Стандартизированные пайплайны предобработки}

В условиях роста числа открытых EEG-датасетов и задач межсубъектного переноса (transfer learning) особое значение приобретает стандартизация ранних этапов обработки. В работе Bigdely-Shamlo и соавт. предложен пайплайн PREP (Preprocessing Pipeline), ориентированный на автоматизированное робастное реферирование и идентификацию плохих каналов~\cite{10.3389/fninf.2015.00016}. Пайплайн включает автоматическую генерацию отчётов и ориентирован на обеспечение воспроизводимости крупномасштабных исследований.

Современные обзоры подчёркивают отсутствие универсального «лучшего» пайплайна предобработки и указывают на зависимость оптимальной стратегии от задачи, типа данных и уровня шума~\cite{10.1109/access.2024.3360328}. Это особенно актуально для задач когнитивного декодирования, где артефакты могут маскировать или имитировать целевые эффекты.

\subsubsection*{Сегментация и нормализация}

После удаления артефактов данные сегментируются на эпохи (epochs), соответствующие временным интервалам интереса. В задачах, связанных с воображением или воспоминанием, эпохи могут иметь значительную длительность, что требует аккуратного выбора окон анализа. Дополнительно применяется коррекция базовой линии (baseline correction) и нормализация амплитуды для снижения межсубъектной вариативности.

\subsubsection*{Выводы по этапу предобработки}

Предобработка EEG представляет собой многоступенчатый процесс, включающий реферирование, фильтрацию, удаление артефактов, сегментацию и нормализацию. Корректность выполнения этих этапов напрямую влияет на устойчивость пространственных и спектральных признаков, используемых в задачах машинного обучения. В контексте реконструкции визуальных стимулов особое значение имеет обеспечение воспроизводимости пайплайна и минимизация артефактных искажений, способных привести к ложной интерпретации пространственно-спектральных закономерностей сигнала.
